\section{模型准备}

% 本章参考优秀论文 A196 的章节组织，用于统一承接“两问及以上共享”的基础机理。
% 若当前赛题不存在跨问共享的实质模型结构，应删除本章，而不是为了编号完整强制保留。
% 本章不写各问最终优化目标、solver、主结果和验证结论。

\subsection{共享机理与基础关系}
% 只定义后续多问都会复用的对象、坐标系、运动/状态关系、守恒关系、公共参数或几何结构。
% 每个核心关系都应说明来源、现实含义以及它将被哪些问题调用。

% 示例占位：
% 设公共状态向量为 $\mathbf{s}(t)$，则基础演化关系可写为
% \begin{equation}
%     \dot{\mathbf{s}}(t)=F\bigl(\mathbf{s}(t),\boldsymbol{\theta}\bigr).
%     \label{eq:shared-state}
% \end{equation}
% 式~\eqref{eq:shared-state} 只表示共享基础关系；正式项目必须替换为当前题目的真实机理。

\subsection{公共判定关系或约束来源}
% 统一定义后续各问反复使用的事件判定、可行域、几何关系、阈值、边界或评价指标。
% 如果判定关系只服务一个问题，则不要放在“模型准备”，应写回对应问题章节。

% 复杂机理题可在此配置 1--2 张核心机理图，帮助解释变量、边界和判定条件；
% 数据处理本身不属于本章，项目级预处理应使用条件章节 05_data.tex。
